\section{Notation}
\label{sec:notation}
Here, we present the notation we use throughout this paper and the typical variable names.
\begin{table}[h!]
    \centering
    \caption{Notation used throughout this paper}
    \vspace{0.15in}
    \begin{tabular}{rp{0.8\textwidth}}
        $SO(n)$ & Group of rotations in $n$-dimensional space\\
        $O(n)$ & Group of rotations and inversion in $n$-dimensional space\\
        $S^2$ & The 2-sphere, surface defined by $x^2+y^2+z^2=1$\\
        $Y^m_\ell$ & Spherical harmonic function of degree $\ell$ and order $m$\\
        $\Y_\ell$ & The collection of spherical harmonic functions of degree $\ell$ for all orders $m$\\
        $\oplus$ & Denotes a direct sum\\
        $\otimes$ & Denotes a tensor product\\
        $\times$ & Denotes a Cartesian product of spaces and also denotes cross products\\
        $\O$ & Big O notation\\
        $\tosphere$ & Function which takes in spherical harmonic coefficients consisting of single copies of irrep up to some cutoff $L$ and converts it into a spherical signal $f:S^2\to\R$\\
        $\fromsphere$ & Function takes in a spherical signal $f:S^2\to\R$ and converts it to spherical harmonic coefficients consisting of single copies of irrep up to some cutoff $L$
    \end{tabular}
    \label{tab:notation}
\end{table}

\begin{table}[h!]
    \centering
    \caption{Commonly used meanings of symbols}
    \vspace{0.15in}
    \begin{tabular}{rp{0.8\textwidth}}
        $G$ & Denotes a group\\
        $\rho(g)$ & Representation of a group\\
        $A$ & First input space of a constructed bilinearity\\
        $B$ & Second input space of a constructed bilinearity\\
        $C$ & Output space of a constructed bilinearity\\
        $V$ & First input space of a tensor product operation (fixed equivariant bilinearity)\\
        $W$ & Second input space of a tensor product operation (fixed equivariant bilinearity)\\
        $V$ & Output space of a tensor product operation (fixed equivariant bilinearity)\\
        $T$ & Tensor product operation (fixed equivariant bilinearity $V\times W\to Z$)\\
        $\ell$ & Typically used to denote irrep type for $SO(3)$. For spherical signals, used instead to denote spherical harmonic degree which naturally indexes multiplicities of irrep types for VSTP/ISTPs\\
        $c$ & Indexes multiplicities of an irrep type in given space (channel)\\
        $j$ & Used instead of $\ell$ to denote irrep type for VSTP and general ISTPs\\
        $s$ & Denotes irrep type of spherical signal (ie. our signal is a map $f:S^2\to\R^{2s+1}$)\\
        
    \end{tabular}
    \label{tab:symbols}
\end{table}